\chapter{Introduction}
Industrial facilities are increasingly equipped with vibration sensors that generate thousands of measurements per seconds. During a single test run a electrical motor can easily produce thousands of data points per minutes as demonstrated by the dataset used in this thesis. This amount of data is not meant to be stored but monitored in real time so that even an non expert in the domain can detect faults early before a complete machine failure. While there are is a solution to this issue namely principal component analysis, which is the main topic of this thesis, still rely on static full available datasets and is (still) not optimized for streaming data and real time anomaly detection. \\\\
The early detection of mechanical anomalies, such as unnatural vibration development, is a central piece of predictive maintenance and therefore a valuable asset to industry. Unplanned downtime or intensive maintenance are significantly costs to consider while easy monitoring can reduce faults and extend the lifespan of an machine. From a scientific perspective principal component analysis is relevant because its combination with time delay embedding as proposed by \textcite{rakuschek2025visual} produces a powerful visually interpretable "fingerprint" of vibration signals, yet only a offline method was described. Taking this method and transferring it into a streaming context comes along with some difficulties such as how to update the algorithm to eliminate computational bottlenecks and to adapt each part of the algorithm to the new streaming approach. When this is achieved it enables the transition from an academic proof of concept to an deployable application.\\ \\
The current state of research are two separate concepts to be investigated. On one hand \textcite{rakuschek2025visual} combines time delay embeddings and principal component analysis to reveal periodic behavior of vibration signals but under the restriction that the full signal history has to be present and that the covariance matrix and its eignedecomposition has to be recomputed from scratch each time a new data point arrives. This makes the offline inefficient and not suitable for high interval streaming data. Another concept investigated in this thesis is the incremental principal component analysis approach by \textcite{Incremental_PCA} along with subspace tracking from \textcite{Oja1982SimplifiedNM} which is a tool for iteratively approximating eigenvectors without calculating the full eignedecomposition solving one of the many computational bottlenecks discussed in this thesis. However these two concepts have not yet been combined with time delay embedding and sliding window mechanism into a single casted real time capable streaming application.\\ \\
The central research question of this thesis is to explore if an incremental, streaming version if the algorithm proposed by \textcite{rakuschek2025visual} keep the diagnostic quality of the of the offline implementation while reducing computational cost with incremental updates of the covariance matrix and eigenvectors. Additionally this thesis investigates how the principal component analysis parameters embedding dimension, time delay and window size influence the quality of the resulting fingerprint plot. \\ \\
To solve this challenge a platform independent library written in C/C++ is developed, which is oriented on the offline algorithm with the improvement that running statistics (mean,covariance matrix etc.) as well as the dominant eigenvectors via subspace iteration are updated incrementally instead of recomputing them from scratch with each incoming data point. The thesis is structured as follows: Chapter 2 investigates the theoretical foundations of time delay embedding, principal component analysis and sliding window techniques and explains in detail how the original implementation is structured. Chapter 3 justifies the selection of the used tech stack and validates the used datasets. Chapter 4 explains in detail the concrete implementation of the combination of the three individual concepts and reinforces that with mathematical formulas and code snippets. Chapter 5 compares the concrete implementation against a reference computation and evaluates the resulting fingerprint plots on real vibration datasets with different parameters applied on. Finally chapter 6 and 7 summarize the findings and provide an outlook on future possible applications.